\section{Bilateral Positivity}
\label{sec:positivity}

\statusestablished

The positivity predicate \(P\) is the central non-logical primitive of the Scott theory. In the four-valued reconstruction, statements of the form \(P(\varphi)\) carry positive and negative information independently.

\subsection{The property domain}
\label{sec:property-domain}

Scott's D1, D2, and D3 quantify over properties, so the reconstruction must say what the property domain is before those definitions can be lifted. I deliberately do not fix it as a powerset, as a set of intensional functions, or by any comprehension principle. The Scott branch instead treats properties as an arbitrary type \(\Pi\) equipped with

\begin{itemize}
  \item signed exemplification relations \(\mathrm{ex}^{+},\mathrm{ex}^{-}\subseteq W\times D\times\Pi\);
  \item signed positivity predicates \(P^{+},P^{-}\subseteq W\times\Pi\);
  \item a property-negation operation \(\nu:\Pi\to\Pi\);
  \item a distinguished element \(G\in\Pi\).
\end{itemize}

Nothing forces \(\Pi\) to be closed under any operation beyond \(\nu\), and no instance of \(\Pi\) is privileged. Every theorem below is therefore a statement about all such structures. This is also what makes the finite countermodels legitimate rather than approximate: a model whose property domain is the four-element set \(\{G,\nu G,Z,\nu Z\}\) is an instance of the semantics, not a truncation of an intended larger one.

The reading of \(\nu\) as \emph{negation} is likewise not built in. It is constrained only by whatever is assumed explicitly, principally
\[
  \mathrm{ex}^{+}(w,x,\nu\varphi)\iff\mathrm{ex}^{-}(w,x,\varphi)
\]
together with the directional A1 conditions below. In particular, involutivity
\[
  \nu\nu\varphi=\varphi
\]
is a further assumption rather than part of the signature. The mechanized \(T1_T\) and \(T2^+\) theorems do not use it. The substitution arguments in the remainder of this section, and the rigidity-coupling proposition in Section~\ref{sec:collapse}, do use it; I flag this explicitly at each point rather than treating property negation as if it were the object-language connective of Section~\ref{sec:kernel}.

The Fitting comparison in Section~\ref{sec:fitting-obstruction} differs from the Scott branch precisely here. It replaces this abstract domain by a typed intension/extension architecture in which properties have internal structure, and the resulting freedom to choose which extensions are admissible is what generates the obstruction and the repair analyzed there.

\subsection{Directional decomposition of A1}

Scott's A1 is
\[
  \neg P(\varphi)\leftrightarrow P(\neg\varphi).
\]
A split into two biconditional channel schemas would be redundant: if property negation is involutive, then because A1 is a schema over all properties, substituting \(\neg\varphi\) for \(\varphi\) converts either biconditional into the other. I therefore decompose A1 by direction:
\[
  A1_L:\quad \neginfo{P(\varphi)}\Rightarrow\pos{P(\neg\varphi)},
\]
\[
  A1_R:\quad \pos{P(\neg\varphi)}\Rightarrow\neginfo{P(\varphi)}.
\]
Here and below \(\Rightarrow\) remains a metalanguage implication.

\begin{proposition}[Directional independence of A1]
Neither \(A1_L\) nor \(A1_R\) semantically entails the other in the four-valued positivity semantics.
\end{proposition}
\begin{proof}
For \(A1_L\) without \(A1_R\), let \(v(P(\varphi))=\True\) and \(v(P(\neg\varphi))=\Neither\). The antecedents of all \(A1_L\) instances on this property pair lack negative support, while the mirrored \(A1_R\) instance from \(\pos{P(\varphi)}\) to \(\neginfo{P(\neg\varphi)}\) fails. Conversely, let \(v(P(\varphi))=\False\) and \(v(P(\neg\varphi))=\Neither\). Then the \(A1_R\) antecedents lack positive support, whereas \(\neginfo{P(\varphi)}\) holds without \(\pos{P(\neg\varphi)}\), so \(A1_L\) fails.
\end{proof}

When both directions hold, strong bilateral A1 yields
\[
  \pos{P(\neg\varphi)}\Leftrightarrow\neginfo{P(\varphi)}
\]
and, by substitution under involutive property negation,
\[
  \neginfo{P(\neg\varphi)}\Leftrightarrow\pos{P(\varphi)}.
\]
Thus
\[
  v(P(\neg\varphi))=\operatorname{swap}(v(P(\varphi))),
\]
so property negation exchanges \(\True\) and \(\False\) while preserving \(\Both\) and \(\Neither\).

\subsection{A4 and informative negative rigidity}

Scott's A4,
\[
  P(\varphi)\rightarrow\Box P(\varphi),
\]
has the direct positive-support lifting
\[
  \Rp:\quad \pos{P(\varphi)}\Rightarrow\pos{\Box P(\varphi)}.
\]

A previously considered negative clause,
\[
  \neginfo{P(\varphi)}\Rightarrow\neginfo{\Box P(\varphi)},
\]
is not informative on the fixed S5 control frames. By the bilateral modal clause, \(\neginfo{\Box\psi}\) merely requires some accessible world to negatively support \(\psi\); reflexivity makes the current world such a witness whenever \(\neginfo{\psi}\) already holds.

The informative negative-persistence condition is instead
\[
  R^-_{\mathrm{nec}}:\quad
  \neginfo{P(\varphi)}\Rightarrow\pos{\Box\neg P(\varphi)},
\]
equivalently
\[
  \neginfo{P(\varphi)}\Rightarrow\neginfo{\Diamond P(\varphi)}.
\]

\begin{proposition}[A1 couples the rigidity channels]
Globally assuming \(A1_L\), \(A1_R\), and \(\Rp\) entails \(R^-_{\mathrm{nec}}\).
\end{proposition}
\begin{proof}
Suppose \(\neginfo{P(\varphi)}\) at a world \(w\). By \(A1_L\), \(\pos{P(\neg\varphi)}\). Applying \(\Rp\) to \(\neg\varphi\) gives \(\pos{\Box P(\neg\varphi)}\), so every accessible world positively supports \(P(\neg\varphi)\). Global \(A1_R\) therefore makes every accessible world negatively support \(P(\varphi)\). By the bilateral box clause this is exactly \(\pos{\Box\neg P(\varphi)}\).
\end{proof}

This proposition is the first direct evidence that A1 is not merely a static complement axiom. In the presence of A4 it transfers modal persistence from the positive information channel into the negative one.

\subsection{A local obstruction to classical positivity reflection}

For the classical collapse spine it is useful to isolate the minimal support interface
\[
  D1^+:\quad \pos{G(x)}\land\pos{P(\varphi)}\Rightarrow\pos{\varphi(x)}.
\]
This is not yet the final four-valued definition of Godlikeness.

The classical proof uses a local reflection move
\[
  G(x),Z(x)\Rightarrow P(Z).
\]
Its signed analogue would be
\[
  REF^+:\quad \pos{G(x)}\land\pos{Z(x)}\Rightarrow\pos{P(Z)}.
\]
Even strong A1 together with \(D1^+\) does not entail \(REF^+\).

A glut counterassignment takes
\[
  v(Z(x))=\Both,\qquad v(P(Z))=\False,\qquad v(P(\neg Z))=\True.
\]
Strong A1 is respected, and \(D1^+\) requires positive support for \(\neg Z(x)\), equivalently negative support for \(Z(x)\), which the glut already supplies. Yet \(\pos{P(Z)}\) fails.

A distinct gap counterassignment takes
\[
  v(Z(x))=\True,\qquad v(P(Z))=v(P(\neg Z))=\Neither.
\]
Strong A1 preserves the positivity gap, \(D1^+\) is silent, and again \(\pos{P(Z)}\) fails.

These are local signed counterassignments to the proof interface, not yet full higher-order countermodels. They expose two different classical assumptions hidden in the ordinary reductio: consistency of the relevant exemplification and completeness of the relevant positivity information. Both dimensions will therefore be tracked explicitly in the collapse analysis.
